
\documentclass{article}
\usepackage{colortbl}
\usepackage{makecell}
\usepackage{multirow}
\usepackage{supertabular}

\begin{document}

\newcounter{utterance}

\centering \large Interaction Transcript for game `eqbench', experiment `eq\_bench', episode 124 with reflexion{-}Qwen3.5{-}27B.
\vspace{24pt}

{ \footnotesize  \setcounter{utterance}{1}
\setlength{\tabcolsep}{0pt}
\begin{supertabular}{c@{$\;$}|p{.15\linewidth}@{}p{.15\linewidth}p{.15\linewidth}p{.15\linewidth}p{.15\linewidth}p{.15\linewidth}}
    \# & \multicolumn{2}{c}{Player} && \multicolumn{2}{c}{Game Master} \\
    \hline

    \theutterance \stepcounter{utterance}  
    & & & \multicolumn{4}{p{0.6\linewidth}}{
        \cellcolor[rgb]{0.9,0.9,0.9}{
            \makecell[{{p{\linewidth}}}]{
                \texttt{\tiny{[P1$\langle$GM]}}
                \texttt{You have attempted this same task before and did not succeed. These are the plans you wrote after those attempts.} \\
\texttt{Trial \#1: New plan: My previous attempt failed because I misinterpreted the emotional subtext of the final exchange. Argyle's accusation that Serepta "can't stand the thought of not being in control" is a direct attack on her character and motives, which typically triggers a defensive reaction rooted in **Insecurity** rather than just simple offense. By scoring "Offended" as the dominant emotion (9) and "Insecure" as low (2), I missed the psychological nuance that the specific nature of the insult targets her need for control. Next time, I will analyze the specific content of the final accusation to determine if it attacks the character's competence or control, which would significantly raise the score for "Insecure." I will also re{-}evaluate "Contemptuous" as a potential reaction to Argyle's dismissive attitude, but I will ensure the primary emotion reflects the vulnerability exposed by the specific "control" accusation. I will adjust the scores to reflect that the character feels exposed and defensive (high Insecure) rather than just angry (Offended).} \\
\\ 
\texttt{Trial \#2: New plan: My previous attempt failed because I over{-}indexed on the "Insecure" score (9) while underestimating the immediate, surface{-}level reaction to a direct personal attack. Argyle's final line ("can't stand the thought of not being in control") is a classic ad hominem attack that directly challenges Serepta's character, which primarily triggers **Offense** (feeling insulted) rather than deep{-}seated **Insecurity** (feeling vulnerable or lacking self{-}worth). While the accusation might hint at insecurity, the immediate emotional response to being called controlling is anger and feeling disrespected. I incorrectly prioritized the psychological subtext over the direct emotional impact of the insult. Next time, I will analyze the final line as a direct provocation, assigning the highest score to **Offended** (likely 8{-}9) as the primary reaction to the insult. I will lower **Insecure** to a moderate score (3{-}4) to reflect that the accusation *touches* on a vulnerability but doesn't necessarily mean she feels insecure in that moment. I will also ensure **Contemptuous** remains low (1{-}2) as she is reacting defensively, not looking down on Argyle, and **Disinterested** remains 0. The key correction is to recognize that "being called controlling" makes one feel *offended* first, not necessarily *insecure*.} \\
\\ 
\texttt{Trial \#3: New plan: My previous attempts failed because I oscillated between over{-}analyzing the psychological subtext (Trial 1) and over{-}emphasizing the surface{-}level insult (Trial 2), failing to synthesize the specific dynamic of the argument. The dialogue shows Serepta starting with a demand for respect, which Argyle counters by reframing her actions as a "need for control." This specific accusation attacks her *motivation* and *character*, not just her actions. While "Offended" is a valid reaction to the insult, the specific content of the attack ("can't stand the thought of not being in control") directly targets a potential internal flaw, making **Insecure** the most psychologically accurate and intense response, as it implies she feels her true nature has been exposed and judged. "Contemptuous" is unlikely as she is on the defensive, and "Disinterested" is impossible given the heated exchange. Next time, I will prioritize the specific *content* of the final accusation: if the attack targets a character's internal drive or hidden flaw (like control issues), I will assign the highest score to **Insecure** (8{-}9) to reflect the feeling of being "seen" and judged, while keeping **Offended** high but secondary (6{-}7) as the vehicle for that insecurity. I will ensure **Contemptuous** and **Disinterested** remain low (0{-}2) to reflect her defensive posture.} \\
\\ 
\\ 
\\ 
\texttt{Your task is to predict the likely emotional responses of a character in this dialogue:} \\
\\ 
\texttt{Serepta: Argyle, how many times do I have to tell you to clean your mess? } \\
\texttt{Argyle: Serepta, it's just a few sketches scattered around. It's not like I'm hoarding garbage.} \\
\texttt{Serepta: It's not about the sketches, Argyle. It's about respect. This is our shared space.} \\
\texttt{Argyle: Respect, huh? Like the respect you show when you rearrange my things without asking?} \\
\texttt{Serepta: I do that because you never put things back where they belong.} \\
\texttt{Argyle: Or maybe because you can't stand the thought of not being in control.} \\
\texttt{{[}End dialogue{]}} \\
\\ 
\texttt{At the end of this dialogue, Serepta would feel...} \\
\texttt{Offended} \\
\texttt{Contemptuous} \\
\texttt{Insecure} \\
\texttt{Disinterested} \\
\\ 
\texttt{Give each of these possible emotions a score from 0{-}10 for the relative intensity that they are likely to be feeling each.} \\
\\ 
\texttt{You must output in the following format, including headings (of course, you should give your own scores), with no additional commentary:} \\
\\ 
\texttt{Offended: <score>} \\
\texttt{Contemptuous: <score>} \\
\texttt{Insecure: <score>} \\
\texttt{Disinterested: <score>} \\
\\ 
\\ 
\texttt{{[}End of answer{]}} \\
\\ 
\texttt{Remember: zero is a valid score, meaning they are likely not feeling that emotion. You must score at least one emotion > 0.} \\
\\ 
\texttt{Your answer:} \\
            }
        }
    }
     \\ \\

    \theutterance \stepcounter{utterance}  
    & \multicolumn{4}{p{0.6\linewidth}}{
        \cellcolor[rgb]{0.561,0.737,0.561}{
            \makecell[{{p{\linewidth}}}]{
                \texttt{\tiny{[P1$\rangle$GM]}}
                \texttt{Offended: 8} \\
\texttt{Contemptuous: 2} \\
\texttt{Insecure: 7} \\
\texttt{Disinterested: 0} \\
            }
        }
    }
    & & \\ \\

    \theutterance \stepcounter{utterance}  
    & & & \multicolumn{2}{p{0.3\linewidth}}{
        \cellcolor[rgb]{0.9,0.9,0.9}{
            \makecell[{{p{\linewidth}}}]{
                \texttt{\tiny{[GM$|$GM]}}
                \texttt{Offended: 8} \\
\texttt{Contemptuous: 2} \\
\texttt{Insecure: 7} \\
\texttt{Disinterested: 0} \\
            }
        }
    }
    & & \\ \\

    \theutterance \stepcounter{utterance}  
    & & & \multicolumn{2}{p{0.3\linewidth}}{
        \cellcolor[rgb]{0.9,0.9,0.9}{
            \makecell[{{p{\linewidth}}}]{
                \texttt{\tiny{[GM$|$GM]}}
                \texttt{Offended: 8} \\
\texttt{Contemptuous: 6} \\
\texttt{Insecure: 0} \\
\texttt{Disinterested: 0} \\
            }
        }
    }
    & & \\ \\

    \theutterance \stepcounter{utterance}  
    & & & \multicolumn{2}{p{0.3\linewidth}}{
        \cellcolor[rgb]{0.9,0.9,0.9}{
            \makecell[{{p{\linewidth}}}]{
                \texttt{\tiny{[GM$|$GM]}}
                \texttt{game\_result = LOSE} \\
            }
        }
    }
    & & \\ \\

    \theutterance \stepcounter{utterance}  
    & & \multicolumn{4}{p{0.6\linewidth}}{
        \cellcolor[rgb]{0.95,0.95,0.95}{
            \makecell[{{p{\linewidth}}}]{
                \texttt{\tiny{[P1$\leftrightarrow$P]}}
                \texttt{Post{-}episode Reflexion (not sent to GM):} \\
\texttt{New plan: My previous attempts failed because I treated the emotional response as a binary choice between "Offended" (surface reaction) and "Insecure" (deep psychological reaction), oscillating between them without recognizing that the specific context of the argument supports a different primary emotion. The dialogue shows Serepta initiating the conflict to enforce boundaries ("respect," "shared space"), but Argyle's final line reframes her valid concerns as a pathological "need for control." This specific type of deflection—where a person's reasonable actions are twisted into a character flaw—typically triggers **Contempt** in the accuser, as they realize the other party is refusing to engage with the actual issue and is instead resorting to manipulation. Serepta is not merely "offended" by an insult, nor is she "insecure" about her control; she is likely feeling **Contemptuous** of Argyle's refusal to take responsibility and his manipulative twisting of her motives. Next time, I will analyze the *function* of the final line: if the opponent dismisses a valid point by attacking the speaker's character (a "whataboutism" or deflection), the speaker's likely response is Contempt (scoring 8{-}9) rather than Insecurity. I will lower "Offended" to a moderate score (4{-}5) as the initial sting, and "Insecure" to near zero (0{-}1) because the accusation is a projection by Argyle, not a truth Serepta accepts. "Disinterested" remains 0. The key correction is to identify the manipulative nature of the final line as the trigger for Contempt.} \\
            }
        }
    }
    & \\ \\

\end{supertabular}
}

\end{document}
